\documentclass[11pt]{article}
\usepackage[margin=1.05in]{geometry}
\usepackage{amsmath,amssymb,amsthm,mathtools}
\usepackage{booktabs,microtype,xcolor,graphicx}
\usepackage[shortlabels]{enumitem}
\usepackage[colorlinks=true,linkcolor=blue!45!black,citecolor=blue!45!black,urlcolor=blue!45!black]{hyperref}
\usepackage{xr}
% Statement numbers of the merged Paper C, read from its build. Build
% ../paperC/problem131_switching.tex first.
\externaldocument[C-]{../paperC/problem131_switching}
\hypersetup{pdftitle={Sticky priors and the most likely count vector},pdfauthor={Arjun Pemmasani}}
\newtheorem{theorem}{Theorem}[section]
\newtheorem{lemma}[theorem]{Lemma}
\newtheorem{proposition}[theorem]{Proposition}
\newtheorem{corollary}[theorem]{Corollary}
\theoremstyle{definition}
\newtheorem{definition}[theorem]{Definition}
\newtheorem{example}[theorem]{Example}
\theoremstyle{remark}
\newtheorem{remark}[theorem]{Remark}
\newcommand{\Z}{\mathbb Z}
\newcommand{\R}{\mathbb R}
\newcommand{\E}{\mathbb E}
\newcommand{\Prob}{\mathbb P}
\newcommand{\Var}{\operatorname{Var}}
\newcommand{\Cov}{\operatorname{Cov}}
\newcommand{\supp}{\operatorname{supp}}
\newcommand{\ones}{\mathbf 1}
\emergencystretch=1.5em
\title{Sticky priors and the most likely count vector\\[4pt]
\large Working note S3 (moved out of Paper~C, October 2026)}
\author{Arjun Pemmasani}
\date{October 2026}
\begin{document}
\maketitle

\section{Setting}\label{sec:setting}

This note was Section~6 and Appendix~D of an earlier version of Paper~C,
``Face selection for slowly switching Markov chains''~\cite{paperC}. When
Paper~B~\cite{paperB} was merged into Paper~C in October 2026, the application
to sticky priors moved here. The mathematics is unchanged. The statement
numbers of~\cite{paperC} that we quote are those of the merged paper.

We recall what we need from~\cite{paperC}. A chain on $d$ states has the
transition matrix $I+hQ$, where $Q$ is a generator with rates $q_{ij}\ge0$ for
$i\ne j$ and exit rates $q_i=\sum_{j\ne i}q_{ij}$. It runs for $N$ steps, its
first state has the law $\mu$, and $h=T/N$, $L=\log N$ and $\tau=T/L$. The
composition $n$ counts the visits to each state, and $P_N(n)$ is its law. A
face is a nonempty set $F$ of states, its exit rate $\lambda_F$ is minus the
Perron root of $Q$ restricted to $F$, and for an irreducible face
$\alpha^*_F=l_F\circ r_F$ is its quasi-ergodic law, the entrywise product of
the left and right Perron vectors with $l_F\cdot r_F=1$. Put
$\Phi_\tau(F)=\tau\lambda_F+|F|-1$. First-order face selection
(Theorem~\ref*{C-thm:first} of~\cite{paperC}) says that for large $N$ every most
likely composition has a support that minimizes $\Phi_\tau$ among the faces that
a path can use, and that its fractions are within $O(1/L)$ of $\alpha^*_F$
when these faces are irreducible.

For two states with $q_{12}=q_{21}=1$, $S_{a,b}(u)$ is the sum of $u^{s(w)}$
over the words $w$ with $a$ letters $1$ and $b$ letters $2$, where $s(w)$ is
the number of switches of $w$, and $S_N=S_{\lfloor N/2\rfloor,\lceil
N/2\rceil}$. The equation $S_N(u)=1$ has exactly one positive root $u_N$. Put
$p_N=1/(1+u_N)$, $z_N=N(1-p_N)$ and $c_0=\frac12\log(\pi/8)$. Paper~A~\cite{paperA}
studies these numbers at every $N$.

\section{Sticky priors and the most likely count vector}\label{sec:sticky}

In a hidden Markov model the hidden states $X_1,\dots,X_N$ form a Markov chain
on $d$ states. Before any data are seen, the model already says how the $N$
time steps are likely to be shared among the states. This is the law $P_N$ of
the count vector $n=(n_1,\dots,n_d)$, where $n_i$ is the number of time steps in
state $i$ (the composition of~\cite{paperC}). A sticky prior makes
the chain leave its current state rarely, which is the regime of~\cite{paperC}.
We ask which states the most likely count vector uses.

Here is what we use from~\cite{paperC}. Suppose the chain leaves its
state with probability $h=\tau\log N/N$ at each step and then picks the next
state from a kernel $\Pi$. For a set $F$ of states (a face) let $\rho_F$ be
the Perron root of $\Pi$ restricted to $F$. The chain stays in $F$ for all $N$
steps with probability about $N^{-\tau(1-\rho_F)}$, and about $N^{|F|-1}$ count
vectors use exactly the states of $F$. Theorem~\ref*{C-thm:first} of~\cite{paperC} says that for
large $N$ the most likely count vector uses a set $F$ that minimizes
$\Phi_\tau(F)=\tau(1-\rho_F)+|F|-1$, and that its fractions $n/N$ are then
close to the law $\alpha^*_F$, the entrywise product of the left and right
Perron vectors of $\Pi$ on $F$. Besides this we use the description of a tie
between several minimizers (Theorem~\ref*{C-thm:window} of~\cite{paperC}) and the exact identity
for a symmetric pair of states (Proposition~\ref*{C-prop:pair} of~\cite{paperC}), once each.

The answer is as follows. As the stickiness decreases, that is, as $\tau$
grows, the most likely count vector moves from one state to larger sets of
states. Among symmetric kernels only the uniform one jumps from one state
straight to all of them.

Fox et al.~\cite{fox2011} add a self-transition bias $\kappa\ge0$ to the HDP
hidden Markov model (HDP stands for hierarchical Dirichlet process), so that
the hidden chain stays in each state longer. In their weak-limit approximation
with base weights $\bar\beta$ and concentration $\omega$, row $j$ of the
transition matrix $(\tilde\pi_{jl})$ is Dirichlet with parameters
$\omega\bar\beta_1,\dots,\omega\bar\beta_j+\kappa,\dots,\omega\bar\beta_d$. By
neutrality, row $j$ off the diagonal, divided by $1-\tilde\pi_{jj}$, is Dirichlet
with parameters $(\omega\bar\beta_l)_{l\ne j}$ and independent of $\tilde\pi_{jj}$. So
$\kappa$ only sets how sticky the chain is, and the kernel of the switches does
not depend on it. In this section we fix every exit probability $1-\tilde\pi_{jj}$ at
the same value $h=\tau L/N$, the scale of the crossing of the two-state walk, and study the
kernel. Under the prior the exit probabilities are random. By the aggregation
property of the Dirichlet law, $1-\tilde\pi_{jj}$ has the law
$\mathrm{Beta}(\omega(1-\bar\beta_j),\omega\bar\beta_j+\kappa)$, with mean
$\omega(1-\bar\beta_j)/(\omega+\kappa)$, so a mean of order $L/N$ needs
$\kappa\asymp N/L$. Mixing over this law favors single states. A single state then
keeps probability $\mu_j\E[\tilde\pi_{jj}^{N-1}]$, which decays only like a power of
$L$ when $\kappa\asymp N/L$, while at a fixed $h$ of order $L/N$ it decays like a power
of $N$. So the mode of the mixture can differ from the mode studied here, and we do not
treat it. All results of this section are about the law given $\Pi$ and $h$.

So let $\Pi=(\pi_{ij})$ be a stochastic matrix on $[d]$ with zero diagonal,
$\tau>0$ and $h=\tau L/N<1$. The hidden chain has transition matrix
$(1-h)I+h\Pi=I+hQ$ with $Q=\Pi-I$, so
$q_i=1$ and $T=\tau L$ (Example~\ref*{C-ex:models}(iii) of~\cite{paperC}). We assume that
$\pi_{ij}>0$ for all $i\ne j$, which we call (A+), and that the start $\mu$ has
full support. Under the prior, \textup{(A+)} holds almost surely. Let
$\rho_F=\varrho(\Pi_F)$ for the principal submatrix $\Pi_F$ on $F$. Then the
exit rate of $F$ is $\lambda_F=1-\rho_F$, and
$Q_F$ and $\Pi_F$ have the same Perron vectors. A vertex (a single state) has
$\rho_F=0$, and
$\rho_{[d]}=1$. If $2\le|F|\le d-1$, then $\Pi_F$ is irreducible by (A+), and
$0<\rho_F<1$ since $\rho_F$ is at least the least row sum of $\Pi_F$ and the
Perron root of the irreducible $\Pi$ exceeds that of each proper principal
submatrix~\cite{bermanplemmons1994}. So
$\Phi_\tau(F)=\tau(1-\rho_F)+|F|-1$, and we put
\[
 \tau_-=\min_{|F|\ge2}\frac{|F|-1}{\rho_F},\qquad
 \tau_+=\max_{F\ne[d]}\frac{d-|F|}{1-\rho_F}.
\]
We call a global mode of $P_N$, that is, a most likely value of the count
vector before any data are seen, given $\Pi$ and $h$, a \emph{conditional mode}.

\begin{corollary}[support and location of the conditional mode]\label{cor:map}
Assume \textup{(A+)}, let $\mu$ have full support, and fix $\tau>0$. Let
$\mathcal F_\tau$ be the set of faces that minimize $\Phi_\tau$. There is
$N_0(\tau)$ such that for $N\ge N_0(\tau)$ every conditional mode $n^*$
satisfies $\supp n^*\in\mathcal F_\tau$ and $n^*/N=\alpha^*_F+O(1/L)$ with
$F=\supp n^*$. If $\mathcal F_\tau=\{F^*\}$, then
$n^*/N\to\alpha^*_{F^*}=l\circ r$, where $l$ and $r$ are the left and right
Perron vectors of $\Pi_{F^*}$ with $l\cdot r=1$. No reversibility is needed.
\end{corollary}

In words, for large $N$ a conditional mode uses exactly the states of a set that
minimizes $\Phi_\tau$, and it shares the $N$ time steps among them in the
proportions $\alpha^*_F$.

\begin{proof}
By (A+) every face is irreducible and, as $\mu$ has full support,
admissible. So Theorem~\ref*{C-thm:first}(iii) and~(iv) of~\cite{paperC} give the claims, with
$\alpha^*_F=l\circ r$.
\end{proof}

The threshold $N_0(\tau)$ is not uniform in $\tau$. At $\tau=\tau_-$ the
vertices tie with the faces that attain $\tau_-$, and
Theorem~\ref*{C-thm:window}(iii) of~\cite{paperC} with $\tau^*=\tau_-$ and a fixed large $t=t_0$
shows that for large $N$ the conditional mode is not a single state at
$\tau=\tau_-(1-\frac{\log L-2t_0}{2L})<\tau_-$. So $N_0(\tau)\to\infty$ as
$\tau\uparrow\tau_-$.

The next proposition says for which $\tau$ the conditional mode uses one state, all
states, or a set in between.

\begin{proposition}[thresholds]\label{prop:thresholds}
Assume \textup{(A+)} and $d\ge3$.
\begin{enumerate}[(a)]
\item The minimizers of $\Phi_\tau$ are exactly the vertices if and only if
$\tau<\tau_-$, and no vertex minimizes $\Phi_\tau$ if $\tau>\tau_-$. The face
$[d]$ is the unique minimizer if and only if $\tau>\tau_+$, and $[d]$ does not
minimize $\Phi_\tau$ if $\tau<\tau_+$.
\item The thresholds satisfy $1<\tau_-\le d-1\le\tau_+$.
\item The equality $\tau_-=d-1$ holds if and only if $\tau_+=d-1$, and both hold if and only if
\begin{equation}\label{eq:directjump}
 \rho_F\le\frac{|F|-1}{d-1}\qquad\text{for every face } F.
\end{equation}
If \eqref{eq:directjump} fails, then $\tau_-<d-1<\tau_+$, and for
$\tau\in(\tau_-,\tau_+)$ every minimizer $F$ of $\Phi_\tau$ has
$2\le|F|\le d-1$.
\end{enumerate}
Consequently, let $\mu$ have full support, fix $\tau$ and let $N\ge N_0(\tau)$
as in Corollary~\ref{cor:map}. If $\tau<\tau_-$, the prior
modes are the $Ne_i$ with $\mu_i=\max_l\mu_l$. If $\tau>\tau_+$, every
conditional mode has full support. If $\tau_-<\tau<\tau_+$, every prior
mode has a support $F$ with $2\le|F|\le d-1$, and $F=F^*$ if $F^*$ is
the unique minimizer of $\Phi_\tau$.
\end{proposition}

\begin{proof}
We compare every face with a vertex and with $[d]$. The function
$\Phi_\tau(F)=\tau(1-\rho_F)+|F|-1$ is affine in $\tau$, so each comparison
changes sign at one value of $\tau$, and $\tau_-$ and $\tau_+$ are the extreme
ones.

(a) A vertex has $\Phi_\tau=\tau$. For $|F|\ge2$ we have $\rho_F>0$ and
\[
 \Phi_\tau(F)-\tau=|F|-1-\tau\rho_F=\rho_F\Bigl(\frac{|F|-1}{\rho_F}-\tau\Bigr).
\]
So $F$ is worse than a vertex exactly when $\tau<(|F|-1)/\rho_F$. This holds
for every $F$ with $|F|\ge2$ exactly when $\tau<\tau_-$, and if $\tau>\tau_-$
some $F$ is better than every vertex. In the same way $\Phi_\tau([d])=d-1$,
and for $F\ne[d]$ we have $1-\rho_F>0$ and
\[
 \Phi_\tau(F)-(d-1)=\tau(1-\rho_F)-(d-|F|)
 =(1-\rho_F)\Bigl(\tau-\frac{d-|F|}{1-\rho_F}\Bigr).
\]
So $F$ is worse than $[d]$ exactly when $\tau>(d-|F|)/(1-\rho_F)$, which gives
the two claims on $[d]$.

(b) In the minimum that defines $\tau_-$, a face with $2\le|F|\le d-1$ has
$\rho_F<1$ and so $(|F|-1)/\rho_F>1$, and $[d]$ gives the value $d-1\ge2$. So
$\tau_->1$ and $\tau_-\le d-1$. In the maximum that defines $\tau_+$, a vertex
gives the value $d-1$, so $\tau_+\ge d-1$.

(c) By (b), $\tau_-=d-1$ exactly when $(|F|-1)/\rho_F\ge d-1$ for all
$|F|\ge2$, and $\tau_+=d-1$ exactly when $(d-|F|)/(1-\rho_F)\le d-1$ for all
$F\ne[d]$. After clearing the denominators both conditions say
$(d-1)\rho_F\le|F|-1$. This inequality always holds for the vertices and for
$[d]$, so each condition is \eqref{eq:directjump}. If
\eqref{eq:directjump} fails, then neither holds, so $\tau_-<d-1<\tau_+$ by
(b), and (a) excludes the vertices and $[d]$ on $(\tau_-,\tau_+)$.

For the consequence, Corollary~\ref{cor:map} says that the support of a prior
mode minimizes $\Phi_\tau$, and (a) and (c) say which faces can do so. If
$\tau<\tau_-$ the conditional modes are vertex compositions, and
$P_N(Ne_i)=\mu_i(1-h)^{N-1}$ is largest for the largest $\mu_i$. Parts (a) to
(c) are checked in Lean.
\end{proof}

We say that $\Pi$ has a \emph{direct jump} when \eqref{eq:directjump} holds.
Then at first order the conditional mode goes from one state straight to all $d$ states at
$\tau=d-1$. For the uniform kernel $\pi_{ij}=1/(d-1)$ the matrix $\Pi_F$ has
constant row sums, so $\rho_F=(|F|-1)/(d-1)$. Hence
\eqref{eq:directjump} says that every set of states is at least as easy to
leave as under the uniform kernel. That kernel has $Q=(J-dI)/(d-1)$, and all its faces tie at
$\tau=d-1$. This is the degeneracy of the complete graph, the Potts chain of~\cite{paperB}, which~\cite[Section~\ref*{C-sec:complete}]{paperC} treats.
The next corollary shows that a direct jump is special.

\begin{corollary}[only the uniform kernel]\label{cor:uniform-kernel}
Let $d\ge3$ and assume \textup{(A+)}. \textup{(a)} If $\Pi$ is symmetric,
then $\Pi$ has a direct jump if and only if $\pi_{ij}=1/(d-1)$ for all
$i\ne j$. \textup{(b)} If $d=3$, then $\Pi$ has a direct jump if and only if
$\pi_{ij}\pi_{ji}\le\frac14$ for all three pairs. In particular every cyclic
kernel $\pi_{i,i+1}=a$, $\pi_{i,i-1}=1-a$ (indices mod $3$) has a direct
jump.
\end{corollary}

\begin{proof}
We test \eqref{eq:directjump} on the pairs. The matrix $\Pi_{\{i,j\}}$ has
zero diagonal and off-diagonal entries $\pi_{ij}$ and $\pi_{ji}$, so
$\rho_{\{i,j\}}=\sqrt{\pi_{ij}\pi_{ji}}$.

(a) Here $\rho_{\{i,j\}}=\pi_{ij}$, so for a pair \eqref{eq:directjump} reads
$\pi_{ij}\le1/(d-1)$. The $d$ rows sum to $1$ and have $d(d-1)$ off-diagonal
entries, so the average of the $\pi_{ij}$ is $1/(d-1)$. If none of them
exceeds the average, they are all equal to it. So a direct jump forces the
uniform kernel. Conversely the uniform kernel has $\rho_F=(|F|-1)/(d-1)$ for
every face, as we saw above, so it satisfies \eqref{eq:directjump} with
equality.

(b) For $d=3$ the faces are the vertices, the three pairs and $[3]$. For the
vertices and $[3]$, \eqref{eq:directjump} reads $0\le0$ and $1\le1$. For a
pair it reads $\sqrt{\pi_{ij}\pi_{ji}}\le\frac12$. The cyclic kernel has
$\pi_{ij}\pi_{ji}=a(1-a)\le\frac14$ for every pair. Both parts are checked in
Lean.
\end{proof}

So for a symmetric kernel that is not uniform there is a range of $\tau$ in
which the conditional mode uses more than one state but not all of them
(Proposition~\ref{prop:thresholds}(c)). Without symmetry a direct jump is
possible, and for three states we can say how likely it is under the prior.

\begin{proposition}[the direct jump under the prior]\label{prop:directjump}
Let $d=3$ and let the base weights be uniform. In the weak-limit sticky
HDP-HMM the rows of $\Pi$ are then independent, and the first off-diagonal
entry of each row has the law $\mathrm{Beta}(\omega/3,\omega/3)$. Write
$\Prob_\omega$ for this law of $\Pi$ and $D$ for the event that $\Pi$ has a
direct jump. Then \textup{(a)} $\Prob_\omega(D)\to\frac14$ as $\omega\to0$,
and \textup{(b)} $\Prob_\omega(D)=\frac{3\sqrt3}{8\pi\omega}(1+o(1))$ as
$\omega\to\infty$.
\end{proposition}

\begin{proof}
Let $B_1,B_2,B_3$ be independent with law $\mathrm{Beta}(\omega/3,\omega/3)$,
and put $\pi_{12}=B_1$, $\pi_{21}=B_2$ and $\pi_{31}=B_3$, so that the other
off-diagonal entries are $1-B_j$. By Corollary~\ref{cor:uniform-kernel}(b),
$D$ is the event that $B_1B_2$, $(1-B_1)B_3$ and $(1-B_2)(1-B_3)$ are all at
most $\frac14$.

(a) As $\omega\to0$ each $B_j$ tends in law to a fair coin on $\{0,1\}$. So
the law of $(B_1,B_2,B_3)$ tends to the uniform law on the eight points of
$\{0,1\}^3$. At these points the three products take only the values $0$ and
$1$, so none of the points lies on the boundary of $D$ (where some product
equals $\frac14$). Hence $\Prob_\omega(D)$ tends to the limit probability of
$D$, which is the fraction of the eight points at which all three products are
$0$. Only $(1,0,1)$ and $(0,1,0)$ have this property (checked in Lean), so the
limit is $\frac28=\frac14$. Part (b) is proved in Appendix~\ref{app:hdp}.
\end{proof}

In words, a direct jump has probability close to $\frac14$ when the
concentration is small and of order $1/\omega$ when it is large. For (b), for
large $\omega$ the set $D$ is a thin tube around the cyclic kernels, with
cross-section of area about $\frac98t^4$ at the logit coordinate $t$ along the
cyclic kernels (Appendix~\ref{app:hdp}), which gives the order $1/\omega$. So
the uniform tie is not typical. At $\omega=1$, Monte Carlo with $10^6$ draws
gives $\Prob_\omega(D)=0.115$. So at $\omega=1$ about $88$ percent of the kernels drawn from the
prior have a range of $\tau$ in which, at first order, the conditional mode sits on a pair
(Proposition~\ref{prop:thresholds}(c)).

The comparison of a single state with a symmetric pair is exact at every $N$.

\begin{corollary}[an exact link to the two-state walk]\label{cor:sticky-pair}
Let $i\ne j$ satisfy $\pi_{ij}=\pi_{ji}=\pi_0>0$ and $\mu_i=\mu_j>0$. We do not
assume \textup{(A+)}. Then for every $N\ge2$ the balanced compositions are the
most likely ones with support $\{i,j\}$, and they are more likely than $Ne_i$
if and only if $\tau>\tau^{\rm pair}(N):=Nu_N/((\pi_0+u_N)L)$. If also
$\mu_i=\mu_j=\max_l\mu_l$, then for every $\tau$ with
$\tau^{\rm pair}(N)<\tau<N/L$ no conditional mode is a single state. Finally,
\[
 \tau^{\rm pair}(N)=\frac1{\pi_0}\Bigl[1-\frac{\frac12\log L-c_0}{L}
 +O\Bigl(\frac{\log L}{L^2}\Bigr)\Bigr],
\]
which tends to the first-order value $1/\rho_{\{i,j\}}=1/\pi_0$.
\end{corollary}

In words, a single state against a symmetric pair is the two-state walk, so the
value of $\tau$ at which the pair overtakes the state is known at every $N$ and
not only in the limit.

\begin{proof}
We apply Proposition~\ref*{C-prop:pair} of~\cite{paperC}. Every exit rate is $1$, so its
hypotheses hold with $a=\pi_0$ and $q=1$, and its odds are $u=\pi_0h/(1-h)$. It
says that the balanced compositions are the most likely ones on $\{i,j\}$ and
that they beat $Ne_i$ exactly when $u>u_N$. Solving for $h$, this is
$h>u_N/(\pi_0+u_N)$, and with $h=\tau L/N$ it is $\tau>\tau^{\rm pair}(N)$.

For the conditional modes, $P_N(Ne_l)=\mu_l(1-h)^{N-1}\le P_N(Ne_i)$ for every $l$.
So when the pair beats $Ne_i$ it beats every single state. The condition
$\tau<N/L$ is $h<1$, and the range of $\tau$ is not empty as
$u_N/(\pi_0+u_N)<1$.

Finally $Nu_N=z_N/p_N$ and $u_N=O(L/N)$, so
$\tau^{\rm pair}(N)=\frac{z_N}{\pi_0L}(1+O(L/N))$. By
Corollary~\ref*{C-cor:paperA} of~\cite{paperC}, $z_N+\frac12\log z_N=L+c_0+O(1/L)$ and $z_N$ is of
order $L$, so $\log z_N=\log L+O(\log L/L)$ and
$z_N=L-\frac12\log L+c_0+O(\log L/L)$. Dividing by $\pi_0L$ gives the claim.
\end{proof}

So far we looked at the single most likely count vector. The last result
gives the total probability that the chain uses only the states of a set $F$.

\begin{proposition}[face masses]\label{prop:mass}
Assume \textup{(A+)} and $0<h<1$. For every face $F$ with $|F|\ge2$ and
$\Lambda_F=1-h\lambda_F$,
\[
 \frac{\min r_F}{\max r_F}\mu(F)\Lambda_F^{N-1}
 \le\Prob(\supp n\subseteq F)\le
 \frac{\max r_F}{\min r_F}\mu(F)\Lambda_F^{N-1},
\]
and for fixed $\tau$, $\Lambda_F^{N-1}=N^{-\tau\lambda_F}e^{O(L^2/N)}$. Also
$\Prob(\text{one state})=(1-h)^{N-1}\sim N^{-\tau}$ and
$\Prob(\supp n=[d])\to1$ for fixed $\tau$.
\end{proposition}

In words, the chain stays inside $F$ with probability
$\mu(F)N^{-\tau\lambda_F}$ up to a bounded factor, and the probability that it
uses every state tends to $1$.

\begin{proof}
We write the probability with the transition matrix restricted to $F$ and
compare the vector $\ones$ with the Perron vector. Let
$U_F=(1-h)I+h\Pi_F\ge0$. The chain stays in $F$ for all $N$ steps with
probability $\Prob(\supp n\subseteq F)=\mu_F^{\top}U_F^{N-1}\ones$, and
$U_Fr_F=\Lambda_Fr_F$. As $r_F/\max r_F\le\ones\le r_F/\min r_F$ entrywise and
$U_F^{N-1}$ has nonnegative entries, we may replace $\ones$ by these two
multiples of $r_F$. This gives the bounds. For fixed $\tau$,
$(N-1)\log(1-h\lambda_F)=-\tau\lambda_FL+O(L^2/N)$. A vertex composition is a
single path that never switches, which gives $\Prob(\text{one state})$. Every
$F\ne[d]$ has $\lambda_F>0$, so each of the finitely many proper faces has
probability tending to $0$, and $\Prob(\supp n=[d])\to1$. The finite bounds are
checked in Lean.
\end{proof}

\begin{remark}[conditional mode versus mass]\label{rem:mapmass}
Let $d\ge3$. For fixed $\tau<\tau_+$ and large $N$ every conditional mode has a
proper support, while $\Prob(\supp n=[d])\to1$. For $\tau<\tau_-$ the conditional mode uses
one state, an event of probability about $N^{-\tau}$. A face of size $k$ holds
about $N^{k-1}$ compositions, so it can carry most of the mass while each of
them is less likely than a vertex. Table~\ref{tab:sticky} shows this. At
$\tau=0.8$ the conditional mode is a single state, yet $\Prob(\text{one state})=0.0122$. At
$\tau=8$ first
order puts the conditional mode on $\{1,2\}$, but at $N=240$ it has full support. So the
conditional mode can be a poor summary of a sticky prior, and face selection tells us
on which states it sits.
\end{remark}

\begin{table}[t]
\centering\small
\begin{tabular}{@{}rlccc@{}}
\toprule
$\tau$ & conditional mode & one state & $\{1,2\}$ & $[3]$\\
\midrule
0.5 & one state & 0.0643 & 0.465 & 0.379\\
0.8 & one state & 0.0122 & 0.423 & 0.531\\
3 & $(120,120,0)$ & $4.3\times10^{-8}$ & 0.129 & 0.871\\
8 & $(111,111,18)$ & $1.2\times10^{-21}$ & 0.0081 & 0.992\\
\bottomrule
\end{tabular}
\caption{The kernel $\Pi$ with rows $(0,0.9,0.1)$, $(0.9,0,0.1)$,
$(0.5,0.5,0)$, the uniform start and $N=240$, where $\tau_-=10/9$ and
$\tau_+=10$. The conditional mode (from exact dynamic programming, with a tie of
three compositions at $\tau=0.5$ and $0.8$) and the probabilities of the supports.}
\label{tab:sticky}
\end{table}


\appendix
\section{The large-concentration limit of the direct jump}\label{app:hdp}

We prove Proposition~\ref{prop:directjump}(b), with $B_1,B_2,B_3$ independent
$\mathrm{Beta}(a,a)$ as in part (a) and $a=\omega/3\ge1$.

\emph{Idea of the proof.} We pass to the logits $\ell_j$ of the $B_j$. In
these coordinates the three conditions that define $D$ say that three linear
forms $M_1,M_2,M_3$ of $\ell$, whose sum is $0$, are bounded above by
quantities of second order in $\ell$. For large $a$ the law of $\ell$
concentrates near $0$ on the scale $a^{-1/2}$. There the bounds are small, so
on $D$ all three forms are nearly $0$. The points where they vanish form the
line of the cyclic kernels, so $D$ is a thin tube around this line. Let $t$ be
the coordinate along the line. At the point $t$ of the line each of the three
bounds is about $\frac12t^2$, so the section of $D$ at $t$ is close to the
triangle $M_1,M_2,M_3\le\frac12t^2$, which has area $\frac98t^4$. So we expect
$\Prob_\omega(D)\approx a^{3/2}\int t^4e^{-cat^2}dt\approx1/a$, and a Gaussian
integral gives the constant.

Before the proof we set up three things, namely the logits and the form of $D$
in them, two bounds on the function that appears, and coordinates along and
across the cyclic kernels.

\emph{Logits.} Put $\ell_j=\log(B_j/(1-B_j))$, $g(y)=2\log\cosh(y/2)$ and
$g_j=g(\ell_j)$. Since $B_j=e^{\ell_j/2}/(2\cosh(\ell_j/2))$ and
$1-B_j=e^{-\ell_j/2}/(2\cosh(\ell_j/2))$, we have
\[
 \log(2B_j)=\tfrac12(\ell_j-g_j),\qquad
 \log\bigl(2(1-B_j)\bigr)=-\tfrac12(\ell_j+g_j),\qquad
 B_j(1-B_j)=\tfrac14e^{-g_j}.
\]
The function $g$ is even and convex with $g(0)=g'(0)=0$, so $g\ge0$. The
density of $B_j$ is proportional to $(B_j(1-B_j))^{a-1}$, and
$dB_j/d\ell_j=B_j(1-B_j)$. So $\ell_j$ has the density $e^{-ag(\ell_j)}/Z$ with
$Z=4^a\Gamma(a)^2/\Gamma(2a)=2\sqrt\pi\,\Gamma(a)/\Gamma(a+\frac12)$, where the
second form is the duplication formula for $\Gamma$.

By the proof of Proposition~\ref{prop:directjump}, $D$ is the event that
$B_1B_2$, $(1-B_1)B_3$ and $(1-B_2)(1-B_3)$ are all at most $\frac14$. We
multiply each product by $4$ and take logarithms with the display above. For
example the first condition becomes
$\frac12(\ell_1-g_1)+\frac12(\ell_2-g_2)\le0$. So $D$ is the set where
\begin{equation}\label{eq:hdp-constraints}
 M_1\le g_1+g_2,\qquad M_2\le g_1+g_3,\qquad M_3\le g_2+g_3,
\end{equation}
with $M_1=\ell_1+\ell_2$, $M_2=\ell_3-\ell_1$ and
$M_3=-\ell_2-\ell_3=-M_1-M_2$. Let $\psi=g_1+g_2+g_3$ (in this appendix $\psi$
is not the Bessel function of~\cite{paperC}). Then the joint density
of $\ell=(\ell_1,\ell_2,\ell_3)$ is $Z^{-3}e^{-a\psi(\ell)}$, and
\begin{equation}\label{eq:hdp-integral}
 \Prob_\omega(D)=Z^{-3}\int_De^{-a\psi(\ell)}d\ell.
\end{equation}

\emph{Two bounds on $g$.} For real $y$,
\begin{align}
 \frac{y^2}4-\frac{y^4}{96}\le g(y)&\le\frac{y^2}4,\tag{g1}\label{eq:hdp-g1}\\
 g(y)&\ge\frac1{10}\min(y^2,|y|).\tag{g2}\label{eq:hdp-g2}
\end{align}
The first says that $g(y)$ is close to $y^2/4$ near $0$, and we use it near
the cyclic kernels. The second is a rough lower bound that we use far from
$0$. For \eqref{eq:hdp-g1} put $x=y/2$, so that $g(y)=2\log\cosh x$. The upper
bound is $\cosh x\le e^{x^2/2}$, which holds term by term in the two power
series. The lower bound says that
$\chi(x)=\log\cosh x-\frac{x^2}2+\frac{x^4}{12}$ is nonnegative, and this holds
as $\chi(0)=\chi'(0)=0$ and $\chi''=x^2-\tanh^2x\ge0$. For \eqref{eq:hdp-g2},
$g''=\frac12\mathrm{sech}^2(y/2)>0.2$ on $|y|\le2$, so $g(y)\ge0.1y^2$ there.
For $|y|\ge2$ convexity and $g(0)=0$ give $g(y)\ge\frac12|y|g(2)>0.43|y|$.

\emph{Coordinates along and across the cyclic kernels.} A cyclic kernel has
$B_1=1-B_2=B_3$, so its logits are $\ell=tv$ with $v=(1,-1,1)$ and $t\in\R$,
and $M_1=M_2=M_3=0$ there. For general $\ell$ let $t=\langle\ell,v\rangle/3$.
We use $(t,M_1,M_2)$ as coordinates. The map $\ell\mapsto(t,M_1,M_2)$ is
linear with determinant $1$, and its inverse is $\ell=tv+w$ with
$w=\frac13(M_1-M_2,2M_1+M_2,M_1+2M_2)\perp v$. So $t$ measures the position
along the line and $w$ the displacement from it. Since $|v|^2=3$,
$|\ell|^2=3t^2+|w|^2$, and a direct computation gives
$|w|^2=\frac13(M_1^2+M_2^2+M_3^2)=\frac23(M_1^2+M_1M_2+M_2^2)$. Write $|M|$
and $|M|_\infty$ for the Euclidean norm and the maximum norm of $(M_1,M_2)$.
As $\lvert M_1M_2\rvert\le\frac12|M|^2$,
\begin{equation}\label{eq:hdp-w}
 \tfrac13|M|^2\le|w|^2\le|M|^2\le2|M|_\infty^2.
\end{equation}

\begin{proof}[Proof of Proposition~\ref{prop:directjump}(b)]
We compute the integral in \eqref{eq:hdp-integral} in four steps. Step 1 shows
that on $D$ the forms $M_i$ are of second order in $\ell$. Step 2 removes the
part of $D$ away from $0$. Step 3 finds the section of $D$ at $t$, and Step 4
integrates over $t$.

\emph{Step 1: an a priori bound.} On $D$, each right side
in \eqref{eq:hdp-constraints} is at most $\psi$, so $M_1\le\psi$, and
$M_1=-M_2-M_3\ge-2\psi$. The same holds for $M_2$. So by \eqref{eq:hdp-g1},
which gives $\psi\le\frac14|\ell|^2$,
\begin{equation}\label{eq:hdp-apriori}
 |M|_\infty\le2\psi(\ell)\le\tfrac12|\ell|^2=\tfrac12(3t^2+|w|^2).
\end{equation}

\emph{Step 2: localization.} Put $\varepsilon=1/12$ and let $U$ be the set
where $|t|\le\varepsilon$ and $|M|_\infty\le\varepsilon$. We show that the
integral over $D\setminus U$ is exponentially small in $a$. Off $U$ we have
$|t|>\varepsilon$ or $|M|>\varepsilon$, so by \eqref{eq:hdp-w}
$|\ell|\ge\max(\sqrt3|t|,|M|/\sqrt3)\ge\varepsilon/\sqrt3$. Also some
coordinate has $|\ell_j|\ge|\ell|/\sqrt3$. So \eqref{eq:hdp-g2} gives
\[
 \psi\ge\gamma:=\tfrac1{10}\min\bigl(|\ell|^2/3,\,|\ell|/\sqrt3\bigr)
 \ge\gamma_0:=\varepsilon^2/90 .
\]
Hence $a\psi\ge(a-1)\gamma_0+\gamma$, that is,
$e^{-a\psi}\le e^{-(a-1)\gamma_0}e^{-\gamma}$ off $U$. The function
$e^{-\gamma(\ell)}$ is integrable over $\R^3$, so
$\int_{D\setminus U}e^{-a\psi}=O(e^{-\gamma_0a})$.

\emph{Step 3: the section of $D$ at $t$.} We show that the section of
$D\cap U$ at $t$ lies between two triangles of area close to $\frac98t^4$.
First, on $D\cap U$ the bounds \eqref{eq:hdp-apriori} and \eqref{eq:hdp-w} give
\[
 |M|_\infty\le\tfrac32t^2+\tfrac12|w|^2\le\tfrac32t^2+|M|_\infty^2
 \le\tfrac32t^2+\varepsilon|M|_\infty,
\]
so $|M|_\infty\le\frac32t^2/(1-\varepsilon)\le2t^2$. So $D\cap U$ lies in the
set $E=\{|t|\le\varepsilon,\ |M|_\infty\le2t^2\}$.

Next we estimate $g_j$ and $\psi$ on $E$. There $|w|^2\le8t^4$ by
\eqref{eq:hdp-w}, and $\ell_j=\pm t+w_j$ with $|w_j|\le|w|$. So
$|\ell_j|\le2|t|$ and $|\ell_j^2-t^2|\le7|t|^3$. Then \eqref{eq:hdp-g1} gives
$|g_j-\frac14t^2|\le2|t|^3$, and with $|\ell|^2=3t^2+|w|^2$ it bounds $\psi$.
Hence on $E$
\begin{equation}\label{eq:hdp-pair}
 \bigl|g_i+g_j-\tfrac12t^2\bigr|\le4|t|^3\quad(i\ne j),
 \qquad\bigl|\psi-\tfrac34t^2\bigr|\le2t^4.
\end{equation}
So on $E$ each right side of \eqref{eq:hdp-constraints} is $\frac12t^2$ up to
$4|t|^3$. Now let
$\mathcal T(\delta)=\{M:M_1,M_2\le\delta,\ M_1+M_2\ge-\delta\}$. This is the
set where $M_1,M_2,M_3\le\delta$. It is a triangle with corners
$(\delta,\delta)$, $(\delta,-2\delta)$ and $(-2\delta,\delta)$, so it has area
$\frac92\delta^2$ and lies in $\{|M|_\infty\le2\delta\}$. Let $D_t$ be the
section of $D\cap U$ at $t$ and put $\delta_\pm=\frac12t^2(1\pm8|t|)\ge0$,
which are $\frac12t^2\pm4|t|^3$. If $M\in D_t$, then $(t,M)\in E$, and
\eqref{eq:hdp-constraints} and \eqref{eq:hdp-pair} give
$M_1,M_2,M_3\le\delta_+$. Conversely, if $M\in\mathcal T(\delta_-)$, then
$|M|_\infty\le2\delta_-\le\min(t^2,\varepsilon)$, so $(t,M)\in E\cap U$, and
\eqref{eq:hdp-pair} gives \eqref{eq:hdp-constraints}. So
$\mathcal T(\delta_-)\subseteq D_t\subseteq\mathcal T(\delta_+)$, and
\begin{equation}\label{eq:hdp-area}
 \tfrac98t^4(1-8|t|)^2\le|D_t|\le\tfrac98t^4(1+8|t|)^2.
\end{equation}

\emph{Step 4: integration.} Let $J=\int_{D\cap U}e^{-a\psi}$. We integrate
first over the section and then over $t$ (the change of coordinates has
Jacobian $1$). By \eqref{eq:hdp-pair} and \eqref{eq:hdp-area}, $J$ lies
between the integrals over $|t|\le\varepsilon$ of
$\frac98t^4(1\mp8|t|)^2e^{-a(\frac34t^2\pm2t^4)}$. We substitute
$t=y/\sqrt a$. Then these are $a^{-5/2}$ times integrals over
$|y|\le\varepsilon\sqrt a$ whose integrands tend to $\frac98y^4e^{-3y^2/4}$ as
$a\to\infty$ and are at most $\frac92y^4e^{-y^2/2}$, since $8|t|<1$ and
$2y^4/a=2t^2y^2\le2\varepsilon^2y^2$. By dominated convergence and
$\int_{\R}y^4e^{-3y^2/4}dy=\frac34\sqrt\pi(\frac43)^{5/2}$,
\[
 J=a^{-5/2}\bigl(\sqrt{3\pi}+o(1)\bigr).
\]
By Step 2 the same holds for the integral over $D$. It remains to find $Z$. As
$\Gamma$ is log-convex, $\Gamma(a+\frac12)^2\le a\Gamma(a)^2$ and
$\Gamma(a+1)^2\le(a+\frac12)\Gamma(a+\frac12)^2$, so
$Z=2\sqrt{\pi/a}\,(1+O(1/a))$. Then \eqref{eq:hdp-integral} gives
\[
 \Prob_\omega(D)=\frac{a^{3/2}}{8\pi^{3/2}}\cdot a^{-5/2}\sqrt{3\pi}\,(1+o(1))
 =\frac{\sqrt3}{8\pi a}(1+o(1))=\frac{3\sqrt3}{8\pi\omega}(1+o(1)).
 \qedhere
\]
\end{proof}


\section*{Computations}
The checks of this note are in the script \texttt{verify\_sticky\_fisher.py} in the
folder \nolinkurl{kagey-problems/problem131/paperC} of
\url{https://github.com/apemm/Kagey-Problems} (its part on sticky priors), and
the finite steps are in the Lean module \nolinkurl{lean/PaperC/Sticky.lean}.

\begin{thebibliography}{9}\small

\bibitem{bermanplemmons1994}
A.~Berman and R.~J.~Plemmons, \emph{Nonnegative Matrices in the
Mathematical Sciences}, SIAM, Philadelphia, 1994.
\url{https://doi.org/10.1137/1.9781611971262}.

\bibitem{fox2011}
E.~B.~Fox, E.~B.~Sudderth, M.~I.~Jordan and A.~S.~Willsky, \emph{A sticky
HDP-HMM with application to speaker diarization}, Ann.~Appl.~Stat.~\textbf{5}
(2011), 1020--1056. \url{https://doi.org/10.1214/10-AOAS395}.

\bibitem{paperA}
A.~Pemmasani, \emph{Uniform Bessel bounds for magnetization crossings in finite
Ising chains}, preprint, 2026.
\href{https://arxiv.org/abs/2610.06935}{arXiv:2610.06935}.

\bibitem{paperB}
A.~Pemmasani, \emph{Vertex crossings in a symmetric Markov multinomial
model}, preprint, 2026.
\href{https://arxiv.org/abs/2610.09258}{arXiv:2610.09258}.

\bibitem{paperC}
A.~Pemmasani, \emph{Face selection for slowly switching Markov chains}, in
preparation. Draft in the folder \nolinkurl{kagey-problems/problem131/paperC}
of \url{https://github.com/apemm/Kagey-Problems}.

\end{thebibliography}

\end{document}
